\documentclass[10pt,a4paper]{amsart}
\usepackage[margin=19mm]{geometry}
\usepackage{amsmath,amssymb,mathtools}
\usepackage[scr=rsfs,cal=euler]{mathalfa}
\usepackage{xfrac}
\usepackage[unicode,hidelinks,bookmarks=false]{hyperref}
\newcommand{\ie}{i.e.\ }
\newcommand{\cf}{cf.\ }
\newcommand{\resp}{respectively\ }
\newcommand{\Subsection}{\subsection}
\providecommand{\nobreakdash}{\nobreak\hbox{-}}
\setlength{\emergencystretch}{2em}
\multlinegap0pt
\usepackage{xcolor}
\usepackage[normalem]{ulem}
\newtheorem{thm}{Theorem}
\newtheorem{la}{Lemma}
\newtheorem{Proofof}{Proofof}
\DeclareMathOperator{\Alg}{Alg}
\DeclareMathOperator{\BT}{BT}
\DeclareMathOperator{\CL}{CL}
\DeclareMathOperator{\D}{D}
\DeclareMathOperator{\Extraction}{Extraction}
\DeclareMathOperator{\Main}{Main}
\newcommand{\Ord}{\mathrm{O}}
\newcommand{\maru}[1]{{\ooalign{\hfill$\small#1$\hfill\crcr\scalebox{1.3}{$\bigcirc$}}}}
\newcommand{\tume}{\hspace*{0.2em}}
\title{Structural Origins of Cubic Complexity in Pebble Motion}
\author{Tomoki Nakamigawa, Tadashi Sakuma}
\date{Source: arXiv:2503.20550v7 (CC BY 4.0)}
\begin{document}
\maketitle

\noindent\emph{Source and licence.} Structural Origins of Cubic Complexity in Pebble Motion. Version arXiv:2503.20550v7 (CC BY 4.0), distributed by arXiv under the Creative Commons Attribution 4.0 International licence, http://creativecommons.org/licenses/by/4.0/, as recorded in the arXiv licence metadata for this exact version and shown on the abstract page https://arxiv.org/abs/2503.20550v7. Full licence notice: \texttt{licenses/arxiv-2503.20550-CC-BY-4.0.txt}.\par\smallskip

\noindent\textit{Editorial scope.} Counted unit: the article's thm main2 (source0.tex lines 1136--1146). The article's complete original proof of it is delivered with it (source0.tex lines 2120--2169, one \texttt{proof} environment of 50 lines, which the article places away from the statement). No mathematics is written, repaired or re-derived: the statement and the proof are the article's own lines, in the article's own notation, and no retained line is substituted or re-flowed.  Retained uncounted as the source-local context the proof needs: lines 1062--1065; lines 1345--1366; lines 1916--1921.  Nothing was omitted from any retained span, so the package carries no omission record.  No retained line contains a citation, so the wrapper carries no bibliography and no bibliography entry.  No retained line contains a figure environment, an image-inclusion command or drawing code, so no image is delivered and the package declares no asset dependency.  Editorial formatting only, in addition: an AMS \texttt{amsart} preamble that restates the article's own declarations and macros used by the retained lines, namely the environments \texttt{thm}, \texttt{la}, \texttt{Proofof} on separate counters, together with the standard packages those lines need.  The retained environments print in this document as Theorem 1, La 1, Proofof 1 in order of appearance, whereas the article prints the same items with the numbering of its own full document.  The complete native source of the article as distributed in the arXiv e-print for this exact version is bundled unchanged as \texttt{sources/175427/source0.tex}.
\begin{thm}\label{main}
A solution sequence to the \textsc{Pebble Motion Problem on Trees} can be computed 
in time linear in the length of the sequence, which is $\Ord(nN + n^2 \log(1+\min\{n, k\})).$
\end{thm}

Consider a tree $X$ (the board tree) serving as the input board graph of a PMT. 
If the order $N$ of a tree $X$ (the board tree) is considerably larger than the number of pebbles $n = |P|$ 
to be placed on it, solving the PMT on an appropriate subtree $Y$ of $X$ can be more efficient---both 
in computation time and in the number of moves---than solving it directly on $X$. Here, $Y$ is chosen 
so that its order is at least $n + k + 1$, where $k$ is the size of the maximum isthmus of $X$, and 
the size of the maximum isthmus of $Y$ is at most $k$. The solution of the PMT on $(Y, P)$ can 
then be transferred back to the original instance $(X, P)$. 
The following Lemma\tume\ref{feasible-subtree} is used to find such a suitable $Y$ within $X$. 
In our proposed algorithm, situations repeatedly arise in which one must find a subtree $Z$ of $X$ 
that contains a specified subtree $W$, has a prescribed order, and does not contain any isthmus 
larger than the maximum isthmus of $X$, and then solve the PMT on $Z$. In these cases as well, 
this lemma is used extensively. 
\begin{la}\label{feasible-subtree}
Let $T$ be a tree such that the size of a maximum isthmus of $T$ is at most $k$. 
For any subtree $T^{\prime}$ of $T$ such that the size of a maximum isthmus of $T^{\prime}$ is also at most $k$, the following statements hold:
\begin{itemize}
\item For any natural number $d$ satisfying $|V(T^{\prime})| \leq d \leq |V(T)|$, 
we can find, in $\Ord(d)$, a $d$-vertex subtree $T^{\prime\prime}$ of $T$ such that the size of a maximum isthmus 
of $T^{\prime\prime}$ is at most $k$ and that $T^{\prime\prime}$ contains $T^{\prime}$ 
as a subgraph.
\end{itemize}
\end{la}

$X(\ell)$ is a subgraph of $H(\ell)$,
the maximum isthmus size of $H(\ell)$ is at most $k$,
and all vertices of $H(\ell)$ are black.
Move the $k+1$ empty spaces into $V(H(\ell))$
while recording the sequence of moves.
\label{mini-puzzle}

\begin{thm}\label{main2}
Let $(G, P)$ be a given instance of the \textsc{Pebble Motion Problem on General Graphs}.  
Suppose that, in addition to the instance, %a spanning breadth-first search tree $T_G$ of $G$, and 
a shortest cycle through each vertex $v \in V(G)$ (or $\emptyset$ if none exists), 
is provided. 
Let $\theta$ be an arbitrary integer such that $k \leq \theta \leq N - n - 1$.
Then a solution sequence of length 
$\Ord(n\D(G)+\frac{n^2\min\{n,\CL(G)\}}{\theta}+n^2\log(1+\min\{n, \theta\}))$ 
can be computed in time %linear in its length, which is  
$\Ord(|E(G)|+n\D(G)+\frac{n^2\min\{n,\CL(G)\}}{\theta}+n^2\log(1+\min\{n, \theta\})).$ 
\end{thm}

\begin{Proofof}{Theorem\tume\ref{main2}}
First, computing the $2$-connected component decomposition of the input graph $G$, as well as 
constructing a breadth-first search spanning tree $T_G$ of $G$, requires 
$\Ord(|E(G)|)$ time. 
However, this computational cost is purely algorithmic overhead and does not affect the length of 
the output solution sequence in any way.

In Step\tume\maru{1} of the Main Algorithm ($\Alg.\Main((G, P), f, g)$), when finding a subtree $T$ 
of order $n + k + 1$ of $T_G$, we apply Lemma\tume\ref{feasible-subtree} separately to each part of $T_G$ 
corresponding to the $2$-connected component decomposition of $G$. 
That is, for each $2$-connected component $X$ of $G$, it suffices to ensure that if a part of the subtree $Y$ 
of $T_G$ that spans $X$ is included in $T$, then the size of that part does not exceed the maximum isthmus 
length of $Y$.
It is the same in that $T$ is decomposed into $\Ord(k)$ size subtrees and represented as binary tree data. 
The difference appears in the computational complexity evaluation. 
%我々がここで議論する2-連結ボードグラフGの上のアルゴリズムは、input board treeの代わりにGの適当な
%全域木Tを用い、パラメータkをmaxmum isthmus sizeではなくk:=N-nと定義しなおすこと以外は、木の上の
%アルゴリズムと変わるところはない。Tを細かいO(k)のsubtreeに分解し、それらを二分木データとして表現する
%ことも同じである。違いは計算量評価の方に現れる。
If the input board graph is $2$-connected then the puzzle is feasible if and only if the board graph is 
not a cycle and has at least two unoccupied spaces. Therefore, no matter how the tree $T$ is 
constructed, $T$ may inevitably have to possess an isthmus with at least $k+1$ vertices, 
and as a result, some of the subtrees in the decomposition of $T$ may have long isthmuses of size
$k+1$ or more. It should be noted that each subtree can contain at most one isthmus of size 
at least $k+1$. In such cases, to create a feasible sub-puzzle of size $\Ord(k)$
(c.f. Step\tume\ref{mini-puzzle} of $\Alg.\Extraction(T, \BT,g^{\prime},t)$), 
the puzzle must be constructed around a vertex $v$ with degree $3$ or higher located somewhere 
on a cycle of $G$ that includes the isthmus. To achieve this, the $\Ord(k)$ pebbles that make up 
the sub-puzzle must be brought around $v$. This operation requires $\Ord((k)\times \min\{n,\CL(G)\})$
moves per each sub-puzzle.
In our algorithm, since we have to solve these sub-puzzles $\Ord((\frac{n}{k})^2)$ times, 
the cost of the above operations for the entire algorithm can be estimated as 
$\Ord((\frac{n}{k})^2 \times k \times \min\{n,\CL(G)\})=\Ord(\frac{n^2\min\{n,\CL(G)\}}{k})$.
%cycle以外の$2$-connected board graphにおいては、$2$個以上の空白があればfeasibleであるから、
%$G$のspanning subtree $T$をどのように工夫して取ろうとしても、$T$が長いisthmusを含み、その結果
%$T$の分解におけるsubtreesの一部がそのisthmus上のO(N-n)のsub path達になってしまう可能性は
%避けられない。そうした場合、O(N-n)サイズのfeasibleなsub-puzzleを作るには、そのisthmusを含む
%Gの部分サイクル上のどこかにある次数３以上の頂点vを中心にしてpuzzleを構成するしかない。
%そのためには、当該sub-puzzleを構成するO(N-n)個のpebble達を、vの周囲に持ってくる必要がある。
%この操作には、一つのsub-puzzle当たり、O((N-n)n)のmoveの手間が必要となる。
%我々のアルゴリズムにおいては、これらのsub-puzzle達をO((n/(N-n))^2)回解かなければならないから、
%上記の手数のコストはアルゴリズム全体ではO((n/(N-n))^2 * (N-n)n)=O(n^3/(N-n))と見積もられる。

Note that, for a given instance $(G, P)$, the feasibility of the puzzle is guaranteed as long as 
$k$ satisfies $k \leq N - n - 1$. 
Therefore, in the above discussion, the constant $k$ in the proof may be replaced with an arbitrary 
fixed integer $\theta$ satisfying $k \leq \theta \leq N - n - 1$, without affecting the correctness of the proof.

The statement of our theorem can be directly derived from the above fact and Theorem\tume\ref{main}. 
\end{Proofof}

\end{document}
